% META: {"id":"1SPE-PRODSCAL-EX-029","chapitre":"1SPE-PRODUIT-SCALAIRE","type_objet":"exercice","capacites":["1SPE-PRODUIT-SCALAIRE-C3"],"capacites_codes":["C3"],"methodes":["M3"],"parcours":3,"competences":["calculer"],"duree_min":5,"mode_creation":"ex_nihilo","sources_inspiration":[],"fichier_tex":"chapitres/1SPE-PRODUIT-SCALAIRE/exercices/1SPE-PRODSCAL-EX-029.tex","corrige_tex":"chapitres/1SPE-PRODUIT-SCALAIRE/corriges/1SPE-PRODSCAL-CO-029.tex","status":"approved","origin":"generated","status_history":["generated","audited","accepted","approved"],"release_acceptance":"RELEASE_OWNER_BATCH_ACCEPTANCE","acceptance_closure_digest":"sha256:9b3ccf9a81c5520fbb7e03b7b2d3e7bf2057a02834b6d908bf3be5f49f3b553f","acceptance_receipt_digest":"sha256:4fad86f0282e0fa4e0419cca1ad6379ae7af5a280c04a4a90880f90a1c044242"}
% BEGIN-VERIFY
% from sympy import *
% k, m = symbols('k m', real=True)
% # u=(k,3), v=(2,m): u.v = 2k+3m = 0
% assert solve(2*k + 3*m, k) == [-3*m/2]
% END-VERIFY
\begin{exercice}{1SPE-PRODSCAL-EX-029}{3}{10}
Déterminer l'ensemble des couples $(k, m)$ tels que les vecteurs $\vec{u}\begin{pmatrix} k \\ 3 \end{pmatrix}$ et $\vec{v}\begin{pmatrix} 2 \\ m \end{pmatrix}$ soient orthogonaux. Représenter cet ensemble dans le plan $(k, m)$.
\end{exercice}
